\chapter{Fazit und Ausblick}
\label{chap:fazit}

Die vorliegende Studienarbeit behandelt die methodische Entwicklung und prototypische Umsetzung eines sensorgestützten BLDC-Motorcontrollers. Der Hardwareentwurf, bestehend aus einer B6-Wechselrichterbrücke mit IPP034N08N5-Leistungs-MOSFETs und IR2104-Gate-Treibern, wurde sowohl theoretisch dimensioniert als auch experimentell verifiziert. Die auf einem STM32G431RB-Mikrocontroller implementierte Softwarearchitektur realisiert eine zuverlässige 6-Step-Blockkommutierung, deren Phasensteuerung auf der kontinuierlichen Positionsauswertung von drei integrierten Hall-Sensoren basiert. Oszilloskopische Messungen an einem Lochraster-Prototyp bestätigen die korrekte Taktung der Gate-Signale, die Einhaltung der erforderlichen Totzeit sowie die Funktionalität der Synchrongleichrichtung. Die prinzipielle Funktionsfähigkeit des Systems konnte im Leerlauf- und Niederlastbereich nachgewiesen werden, während die durchgeführten thermischen und elektrischen Analysen die Grenzen des aktuellen Aufbaus für den Volllastbetrieb bei Phasenströmen über $30\,\mathrm{A}$ aufzeigen. Die berechnete Verlustleistung von etwa $28{,}6\,\mathrm{W}$ kann durch die vorhandene passive Kühlung nicht ausreichend abgeführt werden, sodass ein sicherer Dauerbetrieb unter diesen Bedingungen nicht gewährleistet ist. Für den Einsatz im Hochstrombereich ist daher eine Weiterentwicklung der Hardware erforderlich, insbesondere durch die Umsetzung auf einer mehrlagigen Leiterplatte mit angepassten Leiterbahnquerschnitten und thermischen Vias zur Reduktion des thermischen Widerstands sowie zur Minimierung parasitärer Induktivitäten. Ergänzend ist der Einsatz einer aktiven Luftkühlung zweckmäßig. Insgesamt stellt die entwickelte Hard- und Softwarebasis eine geeignete Grundlage für die Weiterentwicklung zu einem leistungsfähigen und robusten Motorcontroller dar.